%% Hand-authored circuitikz figure -- NOT generated by maestro2tex, only placed by it
%% via the "order" list in configs/anatop_tb_tia.json. The `sch_` prefix keeps it clear
%% of the generated fig_*/tab_* namespace.
%%
%% Topology transcribed from subckt anatop_tiaamp in
%%   simarchive/anatop_tb_tia/tia_2026-07-28/Interactive.9/1/TiaFbTB/netlist/input.scs
%% (Spectre order M(drain gate source bulk)). Netlist net names are kept where the
%% figure refers to them (net5, net12) so the two can be cross-read.
%%
%% Simplified: the R_SET and R_z resistor strings are drawn as single trimmable
%% resistors -- in the netlist each is a segmented poly ladder with an NMOS shunt
%% switch per bit (M13--M18 on TCBias, M19--M24 on LCBias). Matching dummies
%% (ModgenDummy_*) carry no signal and are not drawn. Cross-branch gate connections
%% are identified by net name rather than drawn as buses.
\begin{figure}[htbp]
  \centering
  \resizebox{\linewidth}{!}{%
\ctikzset{bipoles/length=0.9cm, transistors/scale=0.82, resistors/scale=0.85}
\begin{circuitikz}[
    every node/.append style={font=\footnotesize},
    gatelab/.style={font=\scriptsize, inner sep=1pt},
    devlab/.style={font=\scriptsize, text=black!60, inner sep=1pt},
    ilab/.style={font=\scriptsize, align=center, inner sep=2pt, text=black!60},
    railnode/.style={circle, fill=black, inner sep=1.05pt},
    grp/.style={draw=black!30, dashed, line width=0.5pt, rounded corners=3pt},
  ]

% ---- geometry -------------------------------------------------------
\def\xA{0}    \def\xB{3.2}                 % bias branches
\def\xM{5.6}  \def\xT{7.0}  \def\xP{8.4}   % input pair (minus / tail / plus)
\def\xO{12.6}                              % output stage
\def\xEn{-2.6}                             % power-down device
\def\yvdd{11.6} \def\yvss{0.4}
\def\yPm{10.0}  \def\yIn{8.2} \def\yNc{6.4} \def\yNl{5.0} \def\yNm{3.6}
\def\ytail{9.25}                           % tail node, net1
\def\yout{7.6}                             % vout
\def\ycomp{9.0}                            % R_z / C_c corridor
\def\yg{6.6}                               % net5 -> M12 gate corridor
\def\ytap{7.15}                            % net5 -> compensation tap
\def\xc{10.2}                              % compensation riser

% ---- supply rails ---------------------------------------------------
\draw[thick] (\xEn-1.0,\yvdd) -- (14.2,\yvdd) node[right]{\texttt{vdd!}};
\draw[thick] (\xEn-1.0,\yvss) -- (14.2,\yvss) node[right]{\texttt{vss!}};

% ===== power-down device =============================================
\node[pmos](M25) at (\xEn,\yPm){};
\draw (M25.source) -- (\xEn,\yvdd);
\draw (M25.gate) -- ++(-0.5,0) node[gatelab, left]{\texttt{En}};
\node[devlab, right=2pt] at (M25.east){M25};
\draw (M25.drain) -- (\xEn,8.4) -- (\xA,8.4);

% ===== bias branch A : degenerated branch, sets the current ==========
\node[pmos](M1) at (\xA,\yPm){};
\node[nmos](M6) at (\xA,\yNc){};
\node[nmos](M8) at (\xA,\yNm){};
\draw (M1.source) -- (\xA,\yvdd);
\draw (M1.drain) -- (M6.drain);
\draw (M6.source) -- (M8.drain);
\draw (M8.source) -- (\xA,2.9);
\draw (\xA,2.9) to[vR, l_=$R_{\mathrm{SET}}$] (\xA,\yvss);
\node[ilab, right=3pt, align=left] at (\xA+0.15,1.55){6-bit \texttt{TCBias}\\shorts segments};
% M1 diode connection -> V_BP
\draw (M1.gate) -- ++(-0.4,0) |- (\xA,8.4);
\node[railnode] at (\xA,8.4){};
\node[gatelab, right=1pt, fill=white, inner sep=1.5pt] at (\xA+0.05,8.4){$V_{\mathrm{BP}}$};
\node[gatelab, left] at (M6.gate){$V_{\mathrm{BNC}}$};
\node[gatelab, left] at (M8.gate){$V_{\mathrm{BN}}$};
\node[devlab, right=2pt] at (M1.east){M1};
\node[devlab, right=2pt] at (M6.east){M6};
\node[devlab, right=2pt] at (M8.east){M8};

% ===== bias branch B : diode branch, develops V_BNC and V_BN =========
\node[pmos](M0) at (\xB,\yPm){};
\node[nmos](M7) at (\xB,\yNc){};
\node[nmos](M9) at (\xB,\yNm){};
\draw (M0.source) -- (\xB,\yvdd);
\draw (M0.drain) -- (M7.drain);
\draw (M7.source) -- (M9.drain);
\draw (M9.source) -- (\xB,\yvss);
\node[gatelab, left] at (M0.gate){$V_{\mathrm{BP}}$};
% M7 diode -> V_BNC ; M9 diode -> V_BN
\draw (M7.gate) -- ++(-0.35,0) |- (\xB,7.6);
\node[railnode] at (\xB,7.6){};
\node[gatelab, right=1pt, fill=white, inner sep=1.5pt] at (\xB+0.05,7.6){$V_{\mathrm{BNC}}$};
\draw (M9.gate) -- ++(-0.35,0) |- (\xB,5.0);
\node[railnode] at (\xB,5.0){};
\node[gatelab, right=1pt, fill=white, inner sep=1.5pt] at (\xB+0.05,5.0){$V_{\mathrm{BN}}$};
\node[devlab, right=2pt] at (M0.east){M0};
\node[devlab, right=2pt] at (M7.east){M7};
\node[devlab, right=2pt] at (M9.east){M9};

% ===== first stage : tail, input pair, mirror load ===================
\node[pmos](M2) at (\xT,\yPm){};
\draw (M2.source) -- (\xT,\yvdd);
\node[gatelab, left] at (M2.gate){$V_{\mathrm{BP}}$};
\node[devlab, right=2pt, align=left] at (M2.east){M2\\\tiny$12/2.4\times2$};
\draw (M2.drain) -- (\xT,\ytail);
\draw (\xM,\ytail) -- (\xP,\ytail);
\node[railnode] at (\xT,\ytail){};

\node[pmos](M3) at (\xM,\yIn){};
\node[pmos, xscale=-1](M4) at (\xP,\yIn){};
\draw (M3.source) -- (\xM,\ytail);
\draw (M4.source) -- (\xP,\ytail);
\draw (M3.gate) -- ++(-0.7,0) node[gatelab, left]{\texttt{vinm}};
\draw (M4.gate) -- ++(0.7,0) node[gatelab, right]{\texttt{vinp}};
\node[devlab, anchor=south east, align=right] at (\xM-0.42,\yIn+0.2){M3\\\tiny$5/0.75\times4$};
\node[devlab, anchor=south west, align=left] at (\xP+0.42,\yIn+0.2){M4\\\tiny$5/0.75\times4$};

\node[nmos, xscale=-1](M10) at (\xM,\yNl){};
\node[nmos](M11) at (\xP,\yNl){};
\draw (M3.drain) -- (M10.drain);
\draw (M4.drain) -- (M11.drain);
\draw (M10.source) -- (\xM,\yvss);
\draw (M11.source) -- (\xP,\yvss);
\node[devlab, anchor=east] at (\xM-0.42,\yNl){M10};
\node[devlab, anchor=west] at (\xP+0.42,\yNl){M11};
% mirror: M10 diode-connected at net12, gate shared with M11
\draw (M10.gate) -- (M11.gate);
\draw (\xM,6.1) node[railnode]{} -- (\xM+0.55,6.1) -- (\xM+0.55,\yNl);
\node[gatelab, anchor=south west, inner sep=2pt] at (\xM+0.06,6.13){\texttt{net12}};
\node[railnode] at (\xM+0.55,\yNl){};
\node[railnode] at (\xP,\ytap){};
\node[gatelab, anchor=south west, inner sep=2pt] at (\xP+0.06,\ytap){\texttt{net5}};

% ===== second stage : common-source driver + current-source load =====
\node[pmos](M5) at (\xO,\yPm){};
\node[nmos](M12) at (\xO,\yNl){};
\draw (M5.source) -- (\xO,\yvdd);
\draw (M5.drain) -- (M12.drain);
\draw (M12.source) -- (\xO,\yvss);
\node[gatelab, left] at (M5.gate){$V_{\mathrm{BP}}$};
\node[devlab, right=2pt, align=left] at (M5.east){M5\\\tiny$12.6/0.75\times2$};
\node[devlab, right=2pt, align=left] at (M12.east){M12\\\tiny$14/0.75\times8$};
\node[railnode] at (\xO,\yout){};
\draw (\xO,\yout) -- (13.9,\yout) node[right]{\texttt{vout}};
% M12 gate driven from net5
\draw (M12.gate) -- ++(-0.55,0) |- (\xP,\yg);
\node[railnode] at (\xP,\yg){};

% ===== Miller compensation ===========================================
\draw (\xP,\ytap) -- (\xc,\ytap) -- (\xc,\ycomp);
\draw (\xc,\ycomp) to[vR, l=$R_z$] (11.4,\ycomp);
\draw (11.4,\ycomp) to[C, l=$C_c$] (\xO,\ycomp);
\node[railnode] at (\xO,\ycomp){};
\node[ilab, anchor=north west] at (10.5,\ycomp-0.28){6-bit \texttt{LCBias}};

% ===== group outlines ================================================
\node[ilab, anchor=north] at (1.6,\yvss-0.45) {self-biased reference};
\node[ilab, anchor=north] at (7.0,\yvss-0.45) {first stage};
\node[ilab, anchor=north] at (12.6,\yvss-0.45) {second stage};
\draw[grp] (-1.0,\yvss-0.35) -- (4.3,\yvss-0.35);
\draw[grp] (4.7,\yvss-0.35) -- (9.3,\yvss-0.35);
\draw[grp] (11.3,\yvss-0.35) -- (13.9,\yvss-0.35);

\end{circuitikz}
  }
  \caption[{Simplified core of the transimpedance amplifier \texttt{anatop\_tiaamp}.}]{Simplified core of the transimpedance amplifier \texttt{anatop\_tiaamp}, drawn from its netlist. It is a two-stage Miller-compensated operational amplifier: a PMOS input pair M3/M4 on tail source M2 with NMOS current-mirror load M10/M11, followed by common-source driver M12 with PMOS current-source load M5, compensated by $C_c$ (a \SI{22.1}{\micro\metre} square MIM) in series with the nulling resistor $R_z$ from the first-stage output \texttt{net5} to \texttt{vout}. Device sizes are given as $W/L$ in \si{\micro\metre} times the multiplicity. The second stage is class~A, so its sourcing capability is the fixed current in M5 while its sinking capability is whatever M12 can pull --- the asymmetry reported in Table~\ref{tab:tia-drive}. The self-biased reference on the left sets every bias rail from $R_{\mathrm{SET}}$; M25 pulls $V_{\mathrm{BP}}$ to \texttt{vdd} to power the amplifier down. Both resistor strings are drawn as single trimmable resistors: each is a segmented poly ladder with one NMOS shunt switch per trim bit.}
  \label{fig:tia-core}
\end{figure}
